% Complete proofs for the corrected Kamp IV.4 and IV.5 presentations.

\subsection{Modulo-four irredundancy and rational-cut separation}
\label{sec:annex-p09-kamp-iv4-iv5-separations}

\begin{annexitem}{annex:p9.kamp.iv4.field}{thm:p09-kamp-iv4-field}{Every field and unary-operation case modulo four}
\begin{proof}
The three sets \(R_0,R_2,W\) are nonempty and pairwise
disjoint, and their union is \(\mathbb Z\). Their eight
possible unions are exactly \(\mathcal B_4\): the empty
union; the three individual sets; the three unions
\(R_0\cup R_2=V\), \(R_0\cup W=R_2^c\),
\(R_2\cup W=R_0^c\); and the full union.
Complement replaces a selection of these three pieces by
its remaining pieces, and intersection retains the pieces
selected in both operands. Thus this is a field, with the
three stated atoms. In particular every field member contains
either all of \(W\) or none of \(W\).

The empty input gives empty output for both isolation
operations, since no witness exists. The full input gives
empty output as well, since the required false neighbor on
the far side of any witness is absent. These are the first
two isolation rows of the table.

Consider next \(A=R_r\), for \(r=0\) or \(r=2\).
For each integer \(i\), there is a greatest member
\(j<i\) of \(R_r\). For example, the positive
integer distances \(i-j\) to earlier members form a
nonempty set and have a least member. The integer
\(j-1\) is outside \(R_r\), and no member of
\(R_r\) lies between \(j\) and \(i\), by the
choice of \(j\). These give every clause for past
isolation at \(i\). A least member \(j>i\) exists
in the same way; its successor \(j+1\) is outside
\(R_r\), and the intervening interval contains no
member. These give future isolation. Hence both outputs
are \(\mathbb Z\), even when \(i\) itself belongs
to the input class: the chosen witness is still strict.

For \(A=V\) or \(A=W\), the same nearest-witness
argument applies. Each input is unbounded in both directions
and the integer next to any of its points is in the other
parity class. Thus both isolation outputs again equal
\(\mathbb Z\). This proves the rows for the single
residue classes and for both parity sets.

It remains to compute the two complemented-class rows.
Let \(A=R_r^c\), with \(r=0\) or \(r=2\).
For past isolation, a candidate witness \(j\) must be
in \(A\), while \(j-1\) must be in \(R_r\).
Thus \(j\in R_{r+1}\), and \(j+1\) belongs to
\(A\). To avoid an intervening input point between
\(j\) and the current point \(i>j\), it is necessary
that \(i=j+1\). It is also sufficient: the intervening
integer interval is then empty. Consequently the past
isolation output is \(R_{r+2}\). For future isolation,
the far-neighbor condition is \(j+1\in R_r\), so
\(j\in R_{r-1}\). The integer \(j-1\) belongs
to \(A\). Avoiding it as an intervening point forces
\(i=j-1\), and that choice is sufficient. This output
is \(R_{r-2}=R_{r+2}\). For \(r=2\) the two
outputs are \(R_0\); for \(r=0\) they are
\(R_2\). Every necessity and converse in the isolation
table has now been verified.

For the ordinary strict past and future, the empty input
again has no witness. Every other member of the field
contains at least one full residue class modulo four.
Such a class is unbounded in both directions: adding or
subtracting sufficiently large multiples of four produces
a member strictly on either side of any specified integer.
Therefore each nonempty field member has both ordinary
past and ordinary future equal to \(\mathbb Z\).
This proves the two remaining columns, including the full
input and all complements.

Finally take a finite formula whose occurring input letters
all receive values in \(\mathcal B_4\). Its truth set
belongs to that field by structural induction. The atomic
case is the valuation assumption. Negation uses field
complement, and conjunction uses field intersection.
Application of past isolation, future isolation, strict
past, or strict future uses the corresponding column of
the complete table just proved. These are every primitive
constructor in \(\mathcal S''\), so the induction is
complete. Finite truth-functional abbreviations preserve
the same field by their Boolean expansion. In particular,
interpreting the single letter \(q_0\) by \(R_2\)
gives the claimed invariant for every candidate one-input
formula
\cite[ch.~IV, thm.~4 proof, pp.~108--109]{Kamp1968}.
\end{proof}
\end{annexitem}

\begin{annexitem}{annex:p9.kamp.iv4}{thm:p09-kamp-iv4}{The two distinct immediate-neighbor obstructions}
\begin{proof}
The immediate-past image of \(R_2\) consists exactly
of integers \(i\) with \(i-1=4k+2\), equivalently
\(i=4k+3\). Thus it is \(R_3\). The immediate-future
image consists exactly of integers with \(i+1=4k+2\),
equivalently \(i=4k+1\), and is \(R_1\).

Every member of \(\mathcal B_4\) contains both one
and three or contains neither, since those points lie
in its single atom \(W\). The set \(R_3\) contains
three and excludes one; \(R_1\) contains one and
excludes three. Neither set belongs to the field.

If a one-input \(\mathcal S''\)-formula expressed
\(\mathsf{PAST}_1\) on all interpretations, evaluate it
at the input \(R_2\). The field lemma puts its truth
set in \(\mathcal B_4\), while the assumed equality
requires the set \(R_3\), a contradiction. If a formula
instead expressed \(\mathsf{FUTURE}_1\), the same
evaluation would require \(R_1\), giving the other
contradiction. This proves both nonexpressibility claims
at the specified basis and domain
\cite[ch.~IV, thm.~4, pp.~108--109]{Kamp1968}.
\end{proof}
\end{annexitem}

\begin{annexitem}{annex:p9.kamp.rational.local.table}{thm:p09-kamp-rational-local-table}{The cut geometry and all sixteen temporal input pairs}
\begin{proof}
The number \(\alpha=\sqrt2\) is irrational: if
\(\alpha=a/b\) in lowest positive integer terms,
then \(a^2=2b^2\) makes \(a\) even, and
substitution then makes \(b\) even, a contradiction.
Comparing squares also gives \(1<\alpha<3/2\).
Every endpoint \(\alpha+n\), \(n\in\mathbb Z\),
is irrational, as is every \(b_n=\alpha+1/n\).
Otherwise subtracting its rational offset would make
\(\alpha=\sqrt2\) rational. The cut \(\alpha\)
itself is also absent from \(\mathbb Q\).

For the periodic set \(Q\), every rational point lies
strictly between two consecutive endpoints
\(\alpha+n,\alpha+n+1\). It is in \(Q\) exactly
when \(n\) is even. Thus one may choose rational
points on both sides of the current point within that
same open interval; the rational neighborhood between
them has constant membership.

Outside \((\alpha,\alpha+1)_{\mathbb Q}\), the sets
\(Q,Q'\) agree. A rational point outside this interval
has a neighborhood avoiding its two irrational endpoints,
and hence has the same local constancy for \(Q'\).
Inside it, the numbers \(b_n\) strictly decrease to
\(\alpha\). For any rational \(x\) in the interval,
the positive real number \(x-\alpha\) lies strictly
between \(1/(n+1)\) and \(1/n\) for a unique
integer \(n\geq1\). Equality would make \(x\)
one of the irrational \(b_n\). Equivalently,
\(b_{n+1}<x<b_n\). The whole rational interval
\((b_{n+1},b_n)_{\mathbb Q}\) is included in \(Q'\)
when \(n\) is odd and excluded when \(n\) is even.
Choose a rational neighborhood of \(x\) within that
interval to obtain local constancy.

This argument handles every rational point, including
points arbitrarily close to \(\alpha\). It does not
assert local constancy at the missing cut or at a real
band endpoint, because those are not domain points.
The unchanged periodic bands arbitrarily far to the
left and right supply points of both \(Q'\) and
its complement in both directions. The same holds for
\(Q\) and its complement. Thus both stated
interpretations have all the required properties.

We prove the operation table for an arbitrary set \(A\)
with these properties. Its four choices are distinct,
since both \(A\) and \(A^c\) are nonempty. The
empty witness set \(X\) gives empty output for both
strict operations. The empty interval set \(Y\)
also gives empty output for every \(X\), since a
strict rational interval between distinct points is
nonempty by density. If \(Y=\mathbb Q\), every
nonempty choice of \(X\) is unbounded in both
directions. At every current point it therefore supplies
both an earlier and a later witness, with the interval
condition automatically satisfied. Thus the last column
is \(\mathbb Q\) except in its empty-witness row.

It remains to treat \(Y=A\) and \(Y=A^c\). Fix
either choice; the local-constancy property applies to
both \(Y\) and its complement. If the current point
\(t\) is outside \(Y\), choose a rational
neighborhood \((u,v)_{\mathbb Q}\) of \(t\)
disjoint from \(Y\). For a proposed past witness
\(s<t\), density supplies a point strictly between
\(\max(s,u)\) and \(t\), and that point is both
intervening and outside \(Y\). For a proposed future
witness \(s>t\), choose a point strictly between
\(t\) and \(\min(s,v)\). It gives the same
contradiction. Hence both outputs are contained in
\(Y\), regardless of \(X\).

If \(t\in Y\) and \(X=Y\) or
\(X=\mathbb Q\), choose a neighborhood of \(t\)
wholly in \(Y\), and take a rational witness strictly
on each side within it. Each witness belongs to
\(X\), and its open interval to \(t\) belongs
to \(Y\). Thus both outputs equal \(Y\) in
these two rows.

The remaining nonempty row is \(X=Y^c\). For a
proposed past witness \(s<t\) in \(Y^c\), local
constancy at \(s\) gives a right-hand neighborhood
outside \(Y\); a rational point in its intersection
with \((s,t)\) contradicts the interval condition.
For a proposed future witness \(s>t\), use the
left-hand neighborhood of \(s\) and a rational point
in its intersection with \((t,s)\). Neither witness
can work. Both outputs are empty in this row. Together
with the already treated empty row this proves every
entry for both choices of \(Y\), and hence all
sixteen input pairs for each of \(S,U\). Strictness
was used throughout; an irrational boundary cannot
serve as an exceptional witness
\cite[ch.~IV, thm.~5 construction and induction,
       pp.~109--112]{Kamp1968}.
\end{proof}
\end{annexitem}

\begin{annexitem}{annex:p9.kamp.iv5.temporal.induction}{thm:p09-kamp-iv5-temporal-induction}{Every Boolean, SINCE and UNTIL induction case}
\begin{proof}
The four candidate pairs are the empty pair, the full
pair, the predicate pair \((Q,Q')\), and its componentwise
complement. Both interpretations have nonempty predicate
and complement, so the four choices are distinct in
each coordinate and correspond unambiguously.

For the atom \(q_0\), the truth-set pair is exactly
\((Q,Q')\). Negation takes complement relative to
\(\mathbb Q\) in each coordinate. It exchanges
the empty and full pairs, and exchanges the predicate
pair with its complement. This proves the negation step.
For conjunction, intersect the two operand pairs
coordinatewise. An empty operand yields the empty pair,
a full operand leaves the other unchanged, two equal
predicate or complement choices retain that choice,
and the intersection of predicate and complement is
empty. These exhaust the Boolean intersection cases
and give a listed pair in every case.

Now suppose \(\eta=S(\xi,\theta)\). The two
induction hypotheses identify the witness operand
\(\xi\) with the same row-choice in each model,
and the interval operand \(\theta\) with the same
column-choice. Apply the complete table of
\cref{thm:p09-kamp-rational-local-table}, once with
\(A=Q\) and once with \(A=Q'\). Every table entry
is the corresponding empty, full, predicate, or
complement choice in both models. Thus the resulting
pair again belongs to the four listed possibilities.
This argument includes empty and full operands, both
current-point colors, matching and opposed witness
colors, and every strict past interval, by the proved
table rather than a restriction on which operands
can arise.

For \(\eta=U(\xi,\theta)\), use the future part
of the same table. It was proved for all the same
sixteen input pairs, using future witnesses and the
left-hand neighborhood of an opposed witness. Its
entries agree with the past table, so the same
coordinate correspondence holds. This supplies the
entire future step as well. Atomic formulas,
negation, conjunction, \(S\), and \(U\) are the
complete \(\mathsf{TL}_1\) grammar, establishing
the pair invariant for every finite formula with
only the allowed input letter.

For a truth-set pair which is empty or full, all
three comparison assertions are immediate. For the
predicate pair, each assertion is exactly its
displayed hypothesis about predicate membership.
For the complemented predicate pair, equal membership
in the predicate gives equal membership in its
complement. This proves the two within-model
comparisons and the comparison between the models.

Since \(1<\alpha<2\), the rational point three
lies strictly between \(\alpha+1\) and
\(\alpha+2\). It is outside \(Q\), and also
outside \(Q'\), since the altered interval lies
below \(\alpha+1\). The cross-model comparison
with \(r=s=3\) therefore proves equality of
the truth values there for every candidate formula
\cite[ch.~IV, thm.~5 proof, pp.~110--112]{Kamp1968}.
\end{proof}
\end{annexitem}

\begin{annexitem}{annex:p9.kamp.iv5}{thm:p09-kamp-iv5}{The corrected first-order witness and the rational separation}
\begin{proof}
All quantifiers in \(\beta\) range over the rational
domain. Its only predicate symbol is \(P_0\), and
its sole free variable is \(t_0\); these properties
also hold for \(\varphi=\neg\beta\). Thus
\(\varphi\) defines a permitted unary rational
tense. We first evaluate the distinguishing formula
completely, and then use the temporal induction.

\paragraph{No accumulating obstruction in the periodic interpretation.}
In fact \(\mathcal M\models\varphi[t_0\mapsto t]\) for
every rational \(t\). Suppose toward a contradiction
that \(y<t\) witnesses \(\beta(t)\). The formula
requires \(y\notin Q\), so \(y\) belongs to
exactly one false unit band
\[
                  \alpha+2k+1<y<\alpha+2k+2
\]
for an integer \(k\). Put
\(l=\alpha+2k+2\) and \(h=\alpha+2k+3\),
the endpoints of the first true band strictly after
\(y\). The nonvacuity clause of \(\beta\)
provides a true point between \(y\) and \(t\),
and every such point must be above \(l\). Hence
\(t>l\). By density choose a rational point
\[
                            l<z<\min(h,t).
\]
It lies in \(Q\) and strictly between \(y,t\),
so the universal clause of \(\beta\) applies to
\(t_2=z\). That clause demands rational
\(p,q\) with
\[
                   y<p<q<z,\qquad p\in Q,
                                      \qquad q\notin Q.
 \tag{1}
\]
But there is no true point between \(y\) and
\(l\). Thus \(p\in Q\) forces \(p>l\),
and then \(l<p<q<z<h\) puts \(q\) in the
same true band. This contradicts (1). No proposed
\(y\) can witness \(\beta(t)\), proving its
falsehood and hence \(\varphi(t)\) for every
rational current point. This also covers a proposed
starting point with no later true point before \(t\):
it fails the separate existential clause at once.

\paragraph{The obstruction at three in the accumulating interpretation.}
Use \(y=1\). Since \(1<\alpha<2\), this point
lies in \((\alpha-1,\alpha)_{\mathbb Q}\),
which is false in both interpretations. It is before
three. The interval \((b_2,b_1)_{\mathbb Q}\)
is true in \(Q'\) and is contained in \((1,3)\).
It therefore supplies the required true point in
\(\beta(3)\); concretely two is such a point
because \(1<\sqrt2<3/2\).

Let now \(z\in Q'\) be any rational point with
\(1<z<3\). All unchanged true bands below the
altered interval end at or below \(\alpha-1<1\),
and the next unchanged true band above it begins
at \(\alpha+2>3\). Therefore
\(\alpha<z<\alpha+1\). Choose an integer
\(j\geq1\) large enough that
\(1/(2j)<z-\alpha\), equivalently \(b_{2j}<z\).
Choose rational points
\[
             p\in(b_{2j+2},b_{2j+1})_{\mathbb Q},
       \qquad q\in(b_{2j+1},b_{2j})_{\mathbb Q}.
 \tag{2}
\]
Both intervals are nonempty. The first is one of
the inserted true bands, while the second is the
intervening false band. Their order gives
\[
                       1<\alpha<p<q<b_{2j}<z.
\]
Thus \(p\in Q'\), \(q\notin Q'\), and
the two points satisfy the required existential
clause for the initially arbitrary \(t_2=z\).
This proves the universal clause for every relevant
true point, including points arbitrarily close to
the cut. With \(t_1=1\), every conjunct of
\(\beta(3)\) holds in \(\mathcal M'\).
Consequently \(\mathcal M'\not\models
\varphi[t_0\mapsto3]\), whereas the preceding paragraph
gives \(\mathcal M\models\varphi[t_0\mapsto3]\).

No endpoint from (2) was used as a witness: each
is irrational and outside the domain. The actual
witnesses \(p,q\) are rational interior points.
There is likewise no quantified witness at the
limit \(\alpha\). The construction only uses
arbitrarily small bands above that missing point,
which is exactly how the rational cut enters the
first-order distinction.

\paragraph{Nonexpressibility.}
Suppose a finite one-input \(\mathsf{TL}_1\)
formula \(\eta\) expressed the unary tense
defined by \(\varphi\) for every interpretation
on \(\mathbb Q\). Apply the assumed equivalence
to the two permitted interpretations \(Q,Q'\)
at the same rational current point three.
The just-proved first-order values are different.
But \cref{thm:p09-kamp-iv5-temporal-induction}
proves equal values there for every \(\eta\)
whose only input letter is \(q_0\). This is a
contradiction. Hence no such temporal expression
exists, even though \(\varphi\) is itself a
first-order unary tense-defining formula in the
source language.

The proof uses the same fixed rational order in
both models and the same current point. Only
the interpretation of the one unary predicate
changes. The domain is dense and has no endpoints,
but it is not order-complete: its irrational cut
is not realized by a rational point. Thus the
result supplies the stated obstruction to extending
the complete-order expressive-completeness theorem
to all linear orders. It establishes the exact
existence claim of Chapter IV, Theorem 5 through
the corrected two-interpretation construction
\cite[ch.~IV, thm.~5, pp.~109--112]{Kamp1968}.
\end{proof}
\end{annexitem}
